\section{Portas e adaptadores ligados por CDI}
\label{sec:hexagonal-cdi}

A arquitetura hexagonal de Alistair Cockburn
(\citetitle{artigo:cockburn-hexagonal}, 2005) chega à mesma regra por
outro caminho. A intenção declarada é permitir que a aplicação seja
\enquote{equally be driven by users, programs, automated test or batch
scripts} e desenvolvida e testada isolada dos dispositivos e bancos que
vai usar. Uma \termo{porta}{\emph{Port}, no vocabulário de
\emph{ports \& adapters}: o ponto onde a aplicação conversa com o mundo
de fora, definido como uma API.} identifica uma conversa com propósito;
um \termo{adaptador}{\emph{Adapter}: converte a API da porta nos sinais
de um dispositivo concreto (HTTP, SQL, fila) e vice-versa.} traduz essa
conversa para uma tecnologia. No Quarkus, o CDI é quem liga uma ponta à
outra.

\paragraph{A porta e o caso de uso.}
Em \texttt{aplicacao.porta}, uma interface; em \texttt{aplicacao}, o
caso de uso que depende dela pelo construtor, sem anotação:

\begin{lstlisting}[language=Java]
package br.com.turma.pedidos.aplicacao.porta;
public interface Pedidos {                 // porta de saída
    Optional<Pedido> porId(PedidoId id);
    void salvar(Pedido pedido);
}

package br.com.turma.pedidos.aplicacao;
public final class FecharPedido {          // caso de uso, Java puro
    private final Pedidos pedidos;
    private final CobrancaGateway cobranca;
    public FecharPedido(Pedidos pedidos, CobrancaGateway cobranca) {
        this.pedidos = pedidos; this.cobranca = cobranca;
    }
    public Pedido executar(PedidoId id) { /* regras em dominio */ }
}
\end{lstlisting}

\paragraph{O adaptador e o producer.}
Em \texttt{adaptador.saida.persistencia}, o adaptador é um bean comum que
implementa a porta. Em \texttt{config}, um producer cria o caso de uso:

\begin{lstlisting}[language=Java]
package br.com.turma.pedidos.adaptador.saida.persistencia;
@ApplicationScoped
class PedidosPanache implements Pedidos { /* usa PedidoRepository */ }

package br.com.turma.pedidos.config;
@ApplicationScoped
class CasosDeUso {
    @Produces @Singleton
    FecharPedido fecharPedido(Pedidos p, CobrancaGateway c) {
        return new FecharPedido(p, c);
    }
}
\end{lstlisting}

\paragraph{Por que \texttt{@Singleton} no producer.}
O guia \citetitle{doc:quarkus-cdi} recomenda \texttt{@ApplicationScoped}
por padrão, a menos que haja um bom motivo para \texttt{@Singleton}. A
diferença é o \termo{client proxy}{Objeto gerado que fica no lugar do bean
injetado e delega para a instância real, criada sob demanda.}: um bean
\texttt{@Singleton} não tem proxy. Aqui há um motivo: \texttt{FecharPedido}
é \texttt{final}, e o guia \citetitle{doc:quarkus-cdi-reference} explica
que classes \texttt{final} não podem ser \enquote{proxiadas}, a não ser que
o Quarkus transforme o bytecode removendo o \texttt{final}. Um caso de uso
imutável e sem estado não ganha nada com o proxy. O preço, segundo o
mesmo guia: um bean \texttt{@Singleton} não pode ser substituído por
\texttt{QuarkusMock}; num \texttt{@QuarkusTest} use \texttt{@InjectMock}
com \texttt{@MockitoConfig(convertScopes = true)}.

\paragraph{Por que isso combina com o build time.}
Toda a ligação é estática e visível no índice: o tipo da porta, o bean
que a implementa e o producer. O guia de CDI diz que exatamente um bean
deve ser atribuível a cada ponto de injeção, senão \enquote{the build
fails} (com \texttt{UnsatisfiedResolutionException} ou
\texttt{AmbiguousResolutionException}). Faltou um adaptador ou sobrou um:
você descobre no build, não na primeira requisição. Note também que nada é
\texttt{private} onde o container precisa entrar: injeção pelo construtor
(ou pelo parâmetro do producer), como recomenda o guia
\citetitle{doc:quarkus-cdi-reference} e como a
seção~\ref{sec:gizmo} explicou em bytecode.

\paragraph{Trade-offs.}
\begin{itemize}
\item O producer é código a mais para manter. Num serviço pequeno,
  anotar \texttt{FecharPedido} com \texttt{@ApplicationScoped} é uma
  troca razoável; registre a decisão.
\item Uma porta por adaptador pode virar burocracia. Crie portas onde
  existe uma segunda implementação real ou provável: o dublê de teste
  (Parte~\ref{parte:testes}) já conta como segunda implementação.
\item Separar domínio e entidade JPA significa mapear um para o outro.
  A próxima seção discute quando isso vale.
\end{itemize}
